\subsection{LightGBM}
\label{subsubsec:eda_lgbm}

LightGBM (\textit{Light Gradient Boosting Machine}) es un framework de gradient boosting desarrollado por Microsoft en 2017 ~\parencite{2017lightgbm}. Este modelo optimiza el proceso de entrenamiento mediante innovaciones en la forma de construir los árboles y muestrear los datos. Mientras que Random Forest construye árboles de forma independiente y paralela, LightGBM los construye secuencialmente, donde cada nuevo árbol se especializa en corregir los errores de los anteriores.

\subsubsection*{Arquitectura y funcionamiento}

El proceso de $boosting$ funciona de la siguiente manera:

\begin{enumerate}
    \item Se entrena un primer árbol simple que intenta clasificar las instancias según las etiquetas.
    \item Se identifican las que este primer árbol clasificó incorrectamente.
    \item Se entrena un segundo árbol enfocado específicamente en esos casos difíciles.
    \item Se combinan ambos árboles para mejorar la precisión global.
    \item El proceso se repite construyendo decenas o cientos de árboles, cada uno corrigiendo los residuos de los anteriores.
\end{enumerate}

\subsubsection*{Innovaciones clave}

A continuación mostramos algunos aspectos por los cuales este modelo resulta interesante de aplicar al reconocimiento de voz clonada por IA:

\begin{itemize}
    \item \textbf{Crecimiento por hoja (\textit{leaf-wise})}: Los árboles crecen expandiendo la hoja que más puede reducir el error en cada paso, lo que produce árboles más profundos en las regiones problemáticas y más eficientes en las fáciles.
    
    \item \textbf{Muestreo basado en gradientes (\textit{GOSS, Gradient-based One-Side Sampling})}: Durante el entrenamiento, LightGBM identifica qué instancias tienen mayor error y las prioriza. Esto permite entrenar con datasets grandes sin procesar toda la información redundante.
    
    \item \textbf{Agrupación de características (\textit{EFB, Exclusive Feature Bundling})}: LightGBM agrupa las características que no suelen tomar valores iguales a 0 para reducir la dimensionalidad.
\end{itemize}

\subsubsection*{Proceso de clasificación}

Cuando una nueva instancia llega para ser clasificada:

\begin{itemize}
    \item Pasa secuencialmente por todos los árboles construidos.
    \item Cada árbol aporta una corrección o ajuste a la predicción anterior.
    \item La suma ponderada de todas las contribuciones produce la clasificación final (probabilidad de pertenecer a una clase).
\end{itemize}

\subsubsection*{Hiperparámetros}

Vamos a explicar los hiperparámetros más conocidos en los modelos LightGBM: 

\begin{itemize}
	\item \textbf{max depth}: Limita la profundidad máxima que pueden alcanzar los árboles durante el entrenamiento.

	\item \textbf{num leaves}: Determina el número máximo de hojas por árbol.
	
	\item \textbf{learning rate}: También conocido como $shrinkage rate$ o eta \(\eta\), este parámetro controla la velocidad a la que se actualizan los pesos del modelo después de cada iteración de $boosting$.
	
	\item \textbf{n estimators}: Determina el número máximo de iteraciones de $boosting$ que se realizarán, lo que equivale al número de árboles que se construirán.
	
	\item \textbf{reg alpha y reg lambda}: Estos parámetros aplican regularización L1 ($lasso$) y L2 ($ridge$) respectivamente para controlar la complejidad del modelo y prevenir el sobreajuste.
	
	\item \textbf{objective}: Define la tarea de aprendizaje que debe realizar el modelo y la función de pérdida que se optimizará durante el entrenamiento.
	
	\item \textbf{boosting type}: Determina el algoritmo específico que LightGBM utilizará para construir los árboles secuencialmente.
\end{itemize}

\subsubsection*{Ventajas específicas para detección de deepfakes}

Ahora veamos qué ventajas presenta LightGBM para detección de $deepfakes$:

\begin{itemize}
    \item \textbf{Alta precisión}: El enfoque secuencial de corrección de errores suele superar a Random Forest en problemas complejos como la detección de deepfakes.
    
    \item \textbf{Eficiencia computacional}: El muestreo inteligente permite entrenar con grandes volúmenes de audios en menos tiempo que otros métodos de $boosting$.
    
    \item \textbf{Manejo de desequilibrios}: Si existen desbalances entre clases, LightGBM incorpora mecanismos para ponderar las muestras adecuadamente.
    
    \item \textbf{Early stopping}: Se puede configurar para que detenga el entrenamiento cuando añadir más árboles no mejora la validación, evitando sobreajuste.
\end{itemize}